\sectionguard
\section{The Documentary Strata}

Everything contested about Guadalupe becomes tractable once the
surviving documents are laid out in the order in which they can
actually be dated, rather than in the order in which the received
narrative places the events they describe. This section does that,
and does nothing else: it establishes what exists, when each witness
enters the record, and what kind of thing each witness is. Whether
the tradition those witnesses transmit is historically true is the
question of \S3 and \S4, and it is deliberately not answered here.

\subsection{The first printed stratum: 1648 and 1649}

Two books, printed in Mexico City thirteen months apart, are the
foundation of every later telling.

The first is Miguel S\'anchez, \work{Imagen de la Virgen Mar\'ia,
Madre de Dios de Guadalupe, milagrosamente aparecida en la ciudad de
M\'exico} (Mexico, 1648), a Spanish work by a diocesan priest and
preacher. It is not a chronicle. It is a theological and rhetorical
elaboration of the apparition as the fulfilment in New Spain of
Apocalypse~12, the woman clothed with the sun; the narrative of the
four encounters is embedded in a devotional and typological argument.
S\'anchez himself says he could not produce documents for the events
and rests on tradition and on papers he does not
identify.\endnote{Miguel S\'anchez, \work{Imagen de la Virgen
Mar\'ia\ldots}, Mexico, 1648. Not consulted directly for this study;
the book's date, character, and place in the sequence are taken from
the eighteenth-century apparitionist compilation \work{Colecci\'on de
obras y opusculos pertenecientes a la milagrosa aparicion de la
bellisima imagen de Nuestra Se\~nora de Guadalupe} (Madrid, 1785),
prologue, read 2026-07-25 in the John Carter Brown Library copy
digitized at
\url{https://archive.org/details/coleccindeobrasy00unkn} (OCR text),
which dates the first printed \term{Relaci\'on} to ``el a\~no de
1648, en que dio \'a luz la Relaci\'on del milagro Miguel S\'anchez''
and describes it as ``mas que Relaci\'on hist\'orica\ldots\ un elogio
ingenioso y florido de la Imagen de Guadalupe.'' On the absence of
documents, see the nineteenth-century statement quoted at \S3.3,
which reports S\'anchez's own admission; this study does not assert
S\'anchez's wording at first hand.}

The second is the Nahuatl tract published by Luis Laso de la Vega,
chaplain and vicar of the Guadalupe hermitage at Tepeyac, in 1649.
Its running title is \work{Huei tlamahui\c{c}oltica}. Its internal
architecture matters, and it can be read off the 1649 imprint itself:
after the licences and a Nahuatl dedicatory letter signed by Laso de
la Vega, the tract opens the apparition narrative under the heading
that has given the narrative its conventional name.

\begin{witnesstext}{Luis Laso de la Vega, \work{Huei
tlamahui\c{c}oltica} (Mexico: Juan Ruiz, 1649), fol.~1r, heading and
opening (Nahuatl, as printed)}
NICAN MOPOHVA, MOTECPANA INQVENIN YANCVICAN HVEITLAMAHVI\c{C}OLTICA
MONEXITI IN\c{C}ENQVIZCA ICHPOCHTLI SANCTA MARIA DIOS YNANTZIN
TO\c{C}IHVAPILLATOCATZIN, IN ONCAN TEPEYACAC MOTENEHVA
GVADALVPE.\par
Acattopa quimottititzino \c{c}e ma\c{c}ehualtzintli itoca Iuan Diego;
Auh \c{c}atepan monexiti initla\c{c}\`o Ixiptlatzin ynixpan yancuican
Obispo D.~Fray Iuan de Sumarraga. Ihuan inixquich tlamahui\c{c}olli
ye quimochihuilia.\par
Ye iuh m\`atlac xihuitl in opehualoc in atl in tepetl M\`exico, yn ye
omoman in mitl, in chimalli, in ye nohuian ontlamatcamani in
ahu\`acan, intepehu\`acan\ldots\ In huel \`iquac inipan Xihuitl mill y
quinientos, y treinta y vno, quiniuh iquezquilhuioc in metztli
Diziembre\ldots\endnote{Read and collated 2026-07-25 against the page
image of the John Carter Brown Library copy digitized at
\url{https://archive.org/details/hueitlamahuiolti00lass} (scan leaf
index~13, bearing the printed folio number~1 and the signature~A).
The heading is
reproduced in the imprint's own orthography, including
\c{c}-cedilla and the printed grave accents; capitalization follows
the printed display types. The passage is quoted as documentary
evidence of the tract's structure and self-description, not as a
devotional or liturgical text. English of the heading, editorially:
``Here is recounted, and set in order, how the wholly perfect Virgin
Saint Mary Mother of God our Queen newly appeared by a great miracle
there at Tepeyacac, called Guadalupe. First she showed herself to a
poor commoner named Juan Diego; and afterwards her precious image
appeared before the new bishop D.~Fray Juan de Sum\'arraga; together
with all the miracles she works.'' The opening sentence dates the
events ``ten years after the city of Mexico was taken\ldots\ in the
year 1531, a few days into the month of December.'' All English
renderings of Nahuatl, Spanish, and Latin in this study are editorial
working translations, visibly commentary, and are not offered as
received translations.}
\end{witnesstext}

The tract is a composite in six identifiable movements: the two
licences of January 1649; a Nahuatl dedicatory address by Laso de la
Vega to the Virgin; the \work{Nican mopohua} narrative; a description
of the sacred image; a second titled section, \work{Nican motecpana},
``here is set in order,'' listing miracles attributed to the image
and adding biographical matter on Juan Diego and his uncle Juan
Bernardino; a closing section opening \nah{Nican tlantica}, ``here
ends the account''; and, after the \term{Laus Deo}, a Nahuatl prayer
to Our Lady of Guadalupe. That last item is a received devotional
text and is deliberately not reproduced anywhere in this
study.\endnote{Structure read 2026-07-25 in the same scan: the
licences of Baltasar Gonz\'alez, S.J. (9 January 1649) and of
Dr Pedro de Barrientos Lomel\'in, provisor and vicar general (11
January 1649), at the head; the headings \work{NICAN MOTECPANA
INIXQVICH TLAMAHVI\c{C}OLLI YEQVIMOCHIHVILIA INILHVICAC
\c{C}IHVAPILLI TOTLA\c{C}ONANTZIN GVADALVPE} and \nah{NICAN tlantica
inittoloca, ini pohualoca in huei tlamahui\c{c}olli} in the OCR text,
with the first collated against its page image; and the final
\nah{TLATLAVHTILIZTLI IC MOTLATLAVHTILITZINO in ilhuicac tlatoca
\c{c}ihuapilli toTla\c{c}onantzin Guadalupe} after the \term{Laus
Deo}. Under \texttt{guidance/editorial.md} a text offered for vocal
prayer may be reproduced only from an identified received witness and
never composed, translated, or adapted by the project or an AI
contributor; this study prints no part of that prayer and offers no
rendering of it. The John Carter Brown volume is itself a composite
binding: following the 1649 tract in the same digitization is a
separate, later bilingual Spanish--Nahuatl \work{Serm\'on de nuestra
gran Reyna\ldots\ Mar\'ia Sant\'issima de Guadalupe}, with its own
pagination beginning again at ``Pag.~1.'' Nothing in this study
treats that sermon as part of the 1649 publication.}

Two facts about this pair are not in dispute between the schools and
should be fixed before anything else is said. First, both books are
of 1648--1649, one hundred and seventeen years after the events they
narrate. Second, nothing printed transmits the narrative before them.
That second point is not an anti-apparitionist thesis; it is conceded
in the most explicitly apparitionist compilation of the eighteenth
century, which opens by admitting that in the first century after the
image appeared ``no printed account of so admirable an event was
published, as it seems,'' and explains the gap
providentially.\endnote{\work{Colecci\'on de obras y opusculos}
(Madrid, 1785), prologue: ``Es verdad que en el primer siglo, que
corri\'o despu\'es de aparecida en M\'exico esta Celestial Copia, no
se imprimi\'o (seg\'un parece) Relaci\'on alguna de suceso tan
admirable, ya fuese por la escasez de las Imprentas, y costoso de las
impresiones, \'o ya por una sabia disposici\'on de la Divina
Providencia.'' Same witness as the preceding note. The compilation
names S\'anchez, Laso de la Vega and Becerra Tanco as the original
printed relations and Francisco de Florencia as the great
seventeenth-century amplifier.}

\subsection{The seventeenth-century amplification}

After 1649 the tradition acquires its standard apologetic
architecture. Luis Becerra Tanco, a diocesan priest and Nahuatl
scholar, published a short account in 1666 and a fuller posthumous
version, \work{Felicidad de M\'exico}, in 1675; he claims to have
seen and read early Nahuatl relations and Indian annals, which he
does not produce.\endnote{Luis Becerra Tanco, \work{Felicidad de
Mexico en el principio, y milagroso origen, que tubo el Santuario de
la Virgen Mar\'ia N.~Se\~nora de Guadalupe\ldots}, second impression,
procured by Antonio de Gama (Mexico: viuda de Bernardo Calder\'on,
1675); catalogued and rights-marked as public domain (Creative
Commons Public Domain Mark~1.0) in the Universidad de Sevilla copy at
\url{https://archive.org/details/A093073}, metadata read 2026-07-25.
The licence page of the 1675 printing is dated 24 June 1675. Becerra
Tanco's claims about lost Nahuatl sources are reported here at second
hand from the 1785 compilation's prologue and from the
nineteenth-century critical literature (\S3.3); his text was not
collated for this study.} Francisco de Florencia, S.J., produced the
large synthesis \work{La estrella del norte de M\'exico} (1688),
which became the standard reference until the nineteenth
century.\endnote{Named in the 1785 compilation's prologue as the
author who ``adelant\'o mucho, para fundar su credibilidad'' and
whose \work{Estrella del Norte de M\'exico} was reprinted with
difficulty because of its bulk. Not consulted directly.}

Between Becerra Tanco's two versions falls the single most important
juridical instrument the apparitionist tradition possesses: the
\work{Informaciones jur\'idicas} taken in 1666 at the instance of the
chapter of the Guadalupe collegiate church, in which aged witnesses
at Cuauhtitl\'an and in Mexico City deposed to what they had heard
from parents and grandparents about Juan Diego and the apparition,
and painters and physicians gave opinions on the image. Their purpose
was to support a petition to Rome for a proper Office and Mass. Two
things about them must be held together: they are formal sworn
testimony, and they are testimony taken one hundred and thirty-five
years after the event, necessarily at second and third hand, by a
body with an interest in the outcome.\endnote{The 1666 proceedings
are described in the 1785 compilation as the ``instrumentos\ldots\
que\ldots\ se formaron el a\~no de 1666 \'a instancias de varios
devotos de la misma Se\~nora'' for the purpose of obtaining a proper
Office and Mass, and are said there to have been supplemented by
fresh diligences after 1750
(\url{https://archive.org/details/coleccindeobrasy00unkn}, read
2026-07-25). They were printed by Fortino Hip\'olito Vera,
\work{Informaciones sobre la milagrosa aparici\'on de la Sant\'isima
Virgen de Guadalupe, recibidas en 1666 y 1723} (Amecameca, 1889);
that edition is identified from the chronology in Edmundo O'Gorman,
\work{Destierro de sombras} (References), Appendix~VII, no.~39, and
was not itself consulted. The characterization of the witnesses'
age and remove is the standard description of the document in both
schools and is used here only at that level.}

\subsection{The sixteenth-century materials claimed for the tradition}

The apparitionist case does not rest on the printed books. It rests
on the claim that a sixteenth-century documentary substratum exists
and that the 1649 tract preserves it. The custodian of the shrine
publishes an explicit inventory of that substratum, grouped as
Indigenous, mestizo, and Spanish documents, and it is worth
reproducing its principal members at the shrine's own level of
assertion, because a reader deserves to see what the strong case
actually consists of before being told what to think of
it.\endnote{Bas\'ilica de Santa Mar\'ia de Guadalupe,
\work{Documentos Hist\'oricos}, \url{https://virgendeguadalupe.org.mx/documentos-historicos/},
read 2026-07-25. The listing below reproduces that page's own
groupings, datings, and attributions. \emph{These are the shrine's
assertions, not this study's findings}; the ceiling column states
what this study is prepared to say. No independent inspection of any
of the manuscripts named was performed, and none of the datings is
adopted here.}

\begin{stratatable}
\work{Nican mopohua} & pr.~1649; assigned by the shrine to Antonio
Valeriano, c.~1556 & Nahuatl apparition narrative in the 1649 tract &
Text verified in the 1649 imprint; the earlier date and the Valeriano
attribution are contested (\S4.4) and are not adopted\\
\work{Nican motecpana} & pr.~1649; assigned by the shrine to Fernando
de Alva Ixtlilx\'ochitl, 1590 & Nahuatl miracle list and Juan Diego
biography in the same tract & Text verified in the 1649 imprint; the
1590 date and attribution are not adopted\\
\term{C\'odice Escalada} (\term{C\'odice 1548}) & claimed 1548;
surfaced 1995, published 1997 & Small parchment drawing with glosses,
a Sahag\'un signature, and a Valeriano glyph & Disputed artefact
(\S4.4); no authentication is asserted here\\
\work{Anales de Juan Bautista}; \work{Anales de Chimalpahin};
various \work{Anales} & 16th--17th c.\ compilations & Nahuatl
annalistic notices of Tepeyac and of an apparition & Existence and
genre reported at the shrine's level only; no reading of any
manuscript is claimed\\
Testaments (Bartolom\'e L\'opez 1537; Cuauhtitl\'an 1559; Verdugo
Quetzalmamalitzin 1563; and others) & 1537--1577 & Bequests naming
the Guadalupe hermitage & Reported at the shrine's level; these
witness \emph{a shrine and its alms}, not the apparition narrative\\
Cervantes de Salazar, \work{Tres di\'alogos latinos} & 1554 &
Humanist description of Mexico City & Reported at the shrine's level;
witnesses a sanctuary at Tepeyac\\
\work{Informaci\'on} of 1556 & Sept.--Oct.~1556 & Archiepiscopal
testimonial inquiry & Read at length in the critical literature
(\S3.2); its bearing is strongly contested\\
Mapa de Uppsala & c.~1556--1562 & Pictorial map of the Valley of
Mexico & Reported; witnesses a church at Tepeyac\\
Viceroy Mart\'in Enr\'iquez to the king & 23 Sept.~1575 & Official
correspondence & Quoted in the critical literature (\S3.2); assigns
the cult a mid-1550s origin\\
S\'anchez, \work{Imagen de la Virgen Mar\'ia} & 1648 & Spanish
theological narrative & Verified as to date and character; first
printed narrative\\
Laso de la Vega, \work{Huei tlamahui\c{c}oltica} & 1649 & Nahuatl
tract & Verified in the imprint; controlling text of the narrative\\
\work{Informaciones jur\'idicas} & 1666 & Sworn depositions at
135 years' remove & Reported at the level of the published edition;
not collated\\
\end{stratatable}

Three observations about this inventory are load-bearing for
everything that follows, and none of them is polemical.

First, the inventory is heterogeneous in kind. A testamentary bequest
to a hermitage, a map showing a church, and a chapter of annals are
not the same sort of evidence as a narrative of four apparitions, and
they do not become so by appearing in the same list. Almost all the
sixteenth-century items that can be dated with confidence establish
the same, real, and genuinely early fact: that a shrine of Our Lady
under the name Guadalupe existed at Tepeyac and attracted alms and
pilgrimage from at least the 1550s. That is a substantial historical
result. It is not the same result as the apparition narrative, and
the distance between the two is exactly the space in which the modern
dispute is conducted.

Second, the items that would close that distance---an early
manuscript of the \work{Nican mopohua}, the \term{C\'odice
Escalada}, the annalistic entries---are precisely the items whose
dating or authenticity is contested. No sixteenth-century manuscript
of the \work{Nican mopohua} is known: the oldest manuscript copy
usually adduced is later, and the text's own first securely dated
appearance is the 1649 print.\endnote{The absence of a
sixteenth-century manuscript witness of the \work{Nican mopohua} is
common ground: the shrine's own presentation dates the composition
by attribution rather than by a surviving dated manuscript, and the
critical literature (O'Gorman, Appendix~I--II; Garc\'ia Icazbalceta,
\S3.3) is emphatic that no such manuscript has been produced. This
study made no search of Mexican manuscript repositories and asserts
only the state of the published question as of 2026-07-25.}

Third, the shrine's inventory and the critical literature are largely
arguing about the same documents. This is not a case where one side
has evidence and the other has none; it is a case where a shared
corpus of undisputed sixteenth-century material about a shrine is
being connected, or not connected, to a narrative first printed in
the seventeenth. What separates the schools is a judgment about
inference, not a quarrel about whether the papers exist.

\subsection{What this section has and has not established}

Established: the 1648 and 1649 printings and their content; the
internal architecture and self-description of the 1649 tract, read in
the imprint; the seventeenth-century apologetic sequence; the
existence and purpose of the 1666 depositions; and the fact,
conceded on both sides, that no printed narrative of the apparition
precedes 1648.

Not established, and not attempted: the date of composition of the
\work{Nican mopohua}; its authorship; the authenticity of the
\term{C\'odice Escalada}; the content of any manuscript not
inspected; and any conclusion whatever about whether the events of
1531 occurred.
